% ==================================================
% Default Style for Code Listings
% ==================================================

% Define custom pastel and subtle colors
\definecolor{codebg}{RGB}{255,250,225}    % pale yellow (change to mint below if desired)
\definecolor{codekw}{RGB}{0,92,197}       % blue keywords
\definecolor{codekw2}{RGB}{163,21,21}     % red-brown keywords
\definecolor{codekw3}{RGB}{127,0,85}      % violet keywords
\definecolor{codestr}{RGB}{0,128,0}       % green strings/return
\definecolor{codecm}{RGB}{120,120,120}    % gray comments
\definecolor{codeid}{RGB}{30,30,30}       % identifiers
\definecolor{codeaccent}{RGB}{255,140,0}  % orange accents
\definecolor{codedelim}{RGB}{180,0,0}     % punctuation

% A lightweight LaTeX language so LaTeX code looks colorful
\lstdefinelanguage{TeXish}{
  sensitive=true,
  morekeywords=[1]{\\begin,\\end,\\usepackage,\\documentclass,\\section,\\subsection,\\emph,\\textbf,\\textit,\\item,\\label,\\ref,\\cite,\\caption,\\includegraphics,\\color,\\lstset,\\lstdefinestyle,\\lstdefinelanguage,\\NewDocumentEnvironment,\\IfBooleanTF,\\IfValueTF,\\centering},
  morekeywords=[2]{\\newcommand,\\renewcommand,\\DeclareMathOperator,\\title,\\author,\\date,\\maketitle,\\includegraphics},
  morekeywords=[3]{\\LaTeX,\\TeX},
  alsoletter=\\,
  morecomment=[l]{\%},
  morestring=[b]",
  morestring=[b]'
}

% Define colors and style for JSON
\colorlet{punct}{red!60!black}
\definecolor{background}{HTML}{EEEEEE}
\definecolor{delim}{RGB}{20,105,176}
\colorlet{numb}{magenta!60!black}

\lstdefinelanguage{json}{
    basicstyle=\normalfont\ttfamily,
    showstringspaces=false,
    breaklines=true,
    %frame=lines,
    % backgroundcolor=\color{background},
    % Literate programming for coloring punctuation and numbers
    literate=
     *{0}{{{\color{numb}0}}}{1}
      {1}{{{\color{numb}1}}}{1}
      {:}{{{\color{punct}{:}}}}{1}
      {,}{{{\color{punct}{,}}}}{1}
      {\{}{{{\color{delim}{\{}}}}{1}
      {\}}{{{\color{delim}{\}}}}}{1}
      {[}{{{\color{delim}{[}}}}{1}
      {]}{{{\color{delim}{]}}}}{1},
}


\lstset{
  basicstyle=\ttfamily\small,                % Monospaced small font
  backgroundcolor=\color{gray!5},            % Light gray background
  frame=single,                              % Draw a box around the code
  breaklines=true,                           % Allow automatic line breaking
  postbreak=\mbox{\textcolor{red}{$\hookrightarrow$}\space}, % Indicate wrapped lines
  numbers=none,                              % Line numbers on the left
  keepspaces=true,                           % Preserve spaces exactly as in source
  columns=flexible,                          % Improves copying of tabular alignment
  % numberstyle=\tiny\color{gray},           % Style of line numbers
  captionpos=b,                              % Caption position (bottom)
  autogobble=true,                           % Trim indentation from source
  showstringspaces=false                     % Don’t mark spaces in strings
  % Add only color-oriented styling below:
  language=TeXish,
  keywordstyle=\color{codekw}\bfseries,             % [1]
  keywordstyle={[2]\color{codekw2}\bfseries},       % [2]
  keywordstyle={[3]\color{codekw3}\bfseries},       % [3]
  stringstyle=\color{codestr},
  commentstyle=\color{codecm}\itshape,
  identifierstyle=\color{codeid},
  emph={\return},
  emphstyle=\color{codestr}\bfseries,
  literate=
    {*}{{\textcolor{codedelim}{*}}}1
    {=}{{\textcolor{codedelim}{=}}}1
    {(}{{\textcolor{codedelim}{(}}}1
    {)}{{\textcolor{codedelim}{)}}}1
    {[}{{\textcolor{codedelim}{[}}}1
    {]}{{\textcolor{codedelim}{]}}}1
    {,}{{\textcolor{codedelim}{,}}}1
    % {"}{{\texttt{"}}}1                      % needed in windows to display double quotes properly 
}


% to style pyhton code in lstinputlisting, define following
% \lstdefinelanguage{PythonCustom}{